% ====================================================================
%  ISI MODUL BAB 1: TEORI BILANGAN
%  Jangan dikompilasi langsung. Kompilasi salah satu dari:
%    teobil_siswa.tex  -> Edisi Siswa
%    teobil_guru.tex   -> Edisi Guru (ada catatan mengajar & pembahasan)
% ====================================================================
\makeatletter
\@ifundefined{ifguru}{\newif\ifguru}{}
\makeatother

\documentclass[11pt,a4paper]{article}
% ====================================================================
%  GAYA BERSAMA MODUL OLIMPIADE (dipakai teobil_isi.tex & kombin_isi.tex)
%  Semua tampilan hitam-putih / abu-abu supaya hemat tinta.
% ====================================================================
\usepackage[bahasai]{babel}
\usepackage{amsmath, amssymb}
\usepackage{geometry}
\geometry{top=2.5cm, bottom=2.5cm, left=2.5cm, right=2.5cm}
\usepackage[most]{tcolorbox}
\usepackage{enumitem}
\usepackage{array}
\usepackage{comment}
\usepackage{tikz}

\setlength{\parindent}{0pt}
\setlength{\parskip}{5pt}

\tcbset{colback=white, colframe=black, colbacktitle=white, coltitle=black,
  fonttitle=\bfseries, boxrule=0.6pt, arc=1mm,
  before upper={\setlength{\parskip}{5pt}}}

% Kegiatan pembuka
\newtcolorbox{selidiki}{enhanced, breakable, colback=black!4, colframe=black!60,
  colbacktitle=black!15, coltitle=black, title={Ayo Selidiki!}}

% Contoh soal (abu-abu tipis, garis tebal di kiri)
\newtcolorbox{contoh}{enhanced, breakable, colback=black!4, colframe=black!55,
  boxrule=0pt, leftrule=3pt, arc=0pt, left=3mm, before skip=6pt, after skip=6pt}

% Soal untuk dicoba langsung (putih, garis putus-putus)
\newtcolorbox{giliranmu}{enhanced, borderline={0.6pt}{0pt}{black, dashed},
  frame hidden, title={Sekarang Giliranmu!},
  attach boxed title to top left={xshift=3mm, yshift=-2mm},
  boxed title style={colframe=white, colback=white}}

% Ringkasan akhir topik
\newtcolorbox{rangkuman}{enhanced, colback=white, colframe=black!60,
  colbacktitle=black!15, coltitle=black, title={Rangkuman}}

% Garis isian
\newcommand{\kosong}{\rule{1.2cm}{0.4pt}}

% Label level: blok abu-abu kecil
\newcommand{\level}[1]{\par\medskip\noindent\colorbox{black!12}{\strut\textbf{#1}}\par\smallskip}

% Pembahasan soal: hanya muncul di Edisi Guru
\newtcolorbox{kotakjawab}{enhanced, breakable, boxrule=0pt, leftrule=2pt, arc=0pt,
  colback=white, colframe=black, left=2mm, top=1mm, bottom=1mm,
  before skip=3pt, after skip=3pt}
\newcommand{\jawab}[2]{\ifguru
  \begin{kotakjawab}\small\textbf{Jawaban: #1.} #2\end{kotakjawab}\fi}

% Catatan mengajar: hanya muncul di Edisi Guru
\ifguru
  \newtcolorbox{catatanguru}{breakable, colback=white, colframe=black!60,
    colbacktitle=black!15, title={Catatan untuk Guru}}
\else
  \excludecomment{catatanguru}
\fi

\newcommand{\C}[2]{\binom{#1}{#2}}


\title{\textbf{Modul Pembinaan Olimpiade Matematika SMP\\Bab 1: Teori Bilangan Dasar}}
\author{\ifguru Edisi Guru\else Bahan Ajar Siswa (Kelas 7, 8, 9)\fi}
\date{}

\begin{document}

\maketitle
\thispagestyle{empty}

\begin{tcolorbox}[title={Halo Calon Sang Juara!}]
Selamat datang di dunia olimpiade matematika! Di sini, kamu tidak dinilai dari seberapa cepat
kamu menghitung atau seberapa banyak rumus yang kamu hafal. Di sini, kamu dinilai dari
\textbf{cara kamu berpikir}. Jangan takut kalau kamu tidak tahu rumusnya, karena banyak soal
olimpiade bisa dipecahkan hanya dengan logika sehari-hari. Selamat mengeksplorasi!
\end{tcolorbox}

\textbf{Cara memakai modul ini.} Setiap topik dimulai dengan kegiatan \textbf{Ayo Selidiki},
lalu dibagi menjadi beberapa langkah kecil. Di setiap langkah ada contoh, lalu kamu langsung
mencoba sendiri di kotak \textbf{Sekarang Giliranmu}. Di akhir topik ada latihan soal dari
Level 1 sampai Level 5.

\ifguru
\begin{catatanguru}
\textbf{Rencana pertemuan (90 menit per pertemuan):} satu topik untuk satu pertemuan, lalu
satu pertemuan terakhir untuk Soal Campuran sebagai tes kecil.

\textbf{Saran umum.} Biarkan siswa mengerjakan Ayo Selidiki sendiri dulu (10--15 menit) sebelum
menjelaskan. Kotak Giliranmu dikerjakan langsung di kelas. Siswa kelas 7 cukup sampai Level 2;
yang cepat paham lanjut ke Level 3--5.

\textbf{Level 4 dan 5} adalah soal uraian. Nilailah dari kelengkapan alasan, bukan hanya
jawaban akhir. Contoh tulisan bukti yang lengkap ada di kotak jawaban.
\end{catatanguru}
\fi

\newpage
% ====================================================================
\section*{Topik 1: Paritas (Logika Ganjil-Genap)}
% ====================================================================

\begin{catatanguru}
Biarkan siswa mencoba teka-teki uang di akhir Ayo Selidiki sampai mereka yakin memang tidak bisa.
Ingatkan bahwa 0 termasuk bilangan genap.
\end{catatanguru}

\begin{selidiki}
Hitung, lalu tulis apakah hasilnya \textbf{ganjil} atau \textbf{genap}.

\begin{center}
\renewcommand{\arraystretch}{1.6}
\begin{tabular}{|l|*{6}{>{\centering\arraybackslash}p{1.3cm}|}}
\hline
Hitungan & $3 + 5$ & $4 + 6$ & $7 + 2$ & $3 \times 5$ & $4 \times 7$ & $2 \times 6$\\ \hline
Ganjil/Genap? & & & & & & \\ \hline
\end{tabular}
\end{center}

Sekarang jawab:
\begin{enumerate}
    \item Kapan hasil penjumlahan dua bilangan menjadi ganjil?
    \item Kapan hasil perkalian dua bilangan menjadi genap?
    \item Kamu jajan dan harus membayar pas \textbf{sepuluh ribu}. Kamu hanya punya lembaran
          lima ribu dan seribu. Bisakah kamu membayar dengan \textbf{tepat 3 lembar}?
\end{enumerate}

Kalau kamu menjawab ``tidak bisa'' pada nomor 3, kamu benar! Lima ribu dan seribu sama-sama
ganjil. Jumlah 3 bilangan ganjil selalu ganjil, padahal sepuluh ribu genap. Sifat ganjil-genap
seperti ini disebut \textbf{paritas}.
\end{selidiki}

\subsection*{Langkah 1: Aturan ganjil-genap}
Ganjil + ganjil = genap, genap + genap = genap, ganjil + genap = ganjil.
Ganjil $\times$ ganjil = ganjil, dan bilangan apa pun dikali genap hasilnya genap.

Satu hal penting: \textbf{bilangan genap tidak mengubah ganjil-genapnya hasil penjumlahan}.
Jadi cukup hitung ada berapa bilangan ganjil.

\begin{contoh}
\textbf{Contoh.} Apakah $1 + 2 + 3 + \dots + 20$ ganjil atau genap?

\textit{Jawab:} Bilangan ganjilnya 1, 3, 5, \dots, 19, ada 10 buah. Jumlah 10 bilangan ganjil
hasilnya genap. Jadi jawabannya \textbf{genap}.
\end{contoh}

\begin{giliranmu}
\begin{enumerate}[label=\alph*.]
    \item $11 + 22 + 33 + 44 + 55$ adalah bilangan \kosong.
    \item $3 \times 5 \times 7 + 2$ adalah bilangan \kosong.
\end{enumerate}
\end{giliranmu}

\subsection*{Langkah 2: Memakai huruf}
Di olimpiade, sering kali bilangannya tidak diketahui dan diganti huruf. Kita tetap bisa
menentukan ganjil-genapnya.

\begin{contoh}
\textbf{Contoh.} Jika $n$ bilangan bulat, apakah $n \times (n+1)$ selalu genap?

\textit{Jawab:} $n$ dan $n+1$ adalah dua bilangan berurutan, jadi salah satunya pasti genap.
Bilangan apa pun dikali genap hasilnya genap. Jadi \textbf{selalu genap}.
\end{contoh}

\begin{giliranmu}
\begin{enumerate}[label=\alph*.]
    \item Jika $x$ ganjil, maka $x^2 + x$ adalah bilangan \kosong.
    \item Jika $a$ genap dan $b$ ganjil, maka $a \times b + b$ adalah bilangan \kosong.
\end{enumerate}
\end{giliranmu}

\subsection*{Langkah 3: Membuktikan sesuatu tidak mungkin}
Paritas sering dipakai untuk menunjukkan bahwa sesuatu \textbf{mustahil}, seperti teka-teki
uang tadi.

\begin{contoh}
\textbf{Contoh.} Bisakah 7 bilangan ganjil dijumlahkan menjadi 50?

\textit{Jawab:} Jumlah 7 bilangan ganjil pasti ganjil, sedangkan 50 genap. Jadi \textbf{tidak bisa}.
\end{contoh}

\begin{giliranmu}
\begin{enumerate}[label=\alph*.]
    \item Bisakah 4 bilangan ganjil dijumlahkan menjadi 21? \kosong
    \item Bisakah jumlah dua bilangan berurutan sama dengan 2026? \kosong
\end{enumerate}
\end{giliranmu}

\begin{rangkuman}
\begin{enumerate}
    \item Ganjil + ganjil = genap. Ganjil + genap = ganjil. Genap + genap = genap.
    \item Ganjil $\times$ ganjil = ganjil. Kalau ada yang genap, hasil kalinya genap.
    \item Jumlah banyak bilangan ganjil jika banyaknya bilangan ganjil di dalamnya ganjil.
    \item Paritas bisa dipakai untuk membuktikan sesuatu tidak mungkin.
\end{enumerate}
\end{rangkuman}

\textit{Kunci Giliranmu:} Langkah 1: ganjil, ganjil. \quad Langkah 2: genap, ganjil. \quad
Langkah 3: tidak bisa, tidak bisa.

\newpage
\subsection*{Latihan Soal}
\textit{Catatan: Kerjakan tanpa harus menghitung hasil akhirnya secara detail!}

\level{Level 1: Pemanasan}
\begin{enumerate}
    \item Tentukan paritas (ganjil/genap) dari $111 + 222 + 333 + 444$.
          \jawab{genap}{Ganjil + genap + ganjil + genap.}
    \item Tentukan paritas dari $13 \times 15 \times 17 \times 19$.
          \jawab{ganjil}{Semua faktornya ganjil.}
    \item Tentukan paritas dari $(2024 \times 2025) + 2026$.
          \jawab{genap}{Genap + genap.}
    \item Tentukan paritas dari $(99 \times 98) + (97 \times 96)$.
          \jawab{genap}{Kedua perkalian memuat bilangan genap.}
    \item Tentukan paritas dari $1 + 2 + 3 + \dots + 10$.
          \jawab{ganjil}{Ada 5 bilangan ganjil. Hasilnya 55.}
\end{enumerate}

\level{Level 2: Eksplorasi}
\begin{enumerate}
    \setcounter{enumi}{5}
    \item Jika $x$ ganjil dan $y$ genap, apakah $x^2 + y$ ganjil atau genap?
          \jawab{ganjil}{$x^2$ ganjil, ditambah genap tetap ganjil.}
    \item Jika $5a$ adalah bilangan genap, apakah $a$ pasti bilangan genap?
          \jawab{ya}{Kalau $a$ ganjil, $5a$ pasti ganjil.}
    \item Jika $a + b$ ganjil, apakah $a \times b$ pasti ganjil atau pasti genap?
          \jawab{pasti genap}{Salah satu dari $a$ dan $b$ harus genap.}
    \item Jika $3m + 7n$ genap, apakah $m + n$ selalu genap? Jelaskan alasanmu.
          \jawab{ya}{Selisih $(3m + 7n) - (m + n) = 2m + 6n$ selalu genap, jadi keduanya sama
          ganjil-genapnya.}
    \item Benar atau salah: ``Jumlah dua bilangan bulat berurutan selalu ganjil.'' Jelaskan.
          \jawab{benar}{Dari dua bilangan berurutan, satu ganjil dan satu genap.}
\end{enumerate}

\level{Level 3: Tantangan OSK}
\begin{enumerate}
    \setcounter{enumi}{10}
    \item Ada 5 kotak yang masing-masing berisi bola. Jumlah seluruh bola adalah 100. Bisakah
          semua kotak berisi bola dalam jumlah ganjil?
          \jawab{tidak bisa}{Jumlah 5 bilangan ganjil pasti ganjil, sedangkan 100 genap.}
    \item \textbf{(Semut di Papan Catur)} Seekor semut berjalan di papan catur. Setiap langkah ia
          pindah dari kotak putih ke hitam atau dari hitam ke putih. Jika ia mulai di kotak putih
          dan melangkah 2025 kali, di warna apa ia berhenti?
          \jawab{hitam}{Setelah langkah ganjil, semut selalu di hitam.}
    \item Di antara angka $1, 2, 3, \dots, 9$ kamu boleh meletakkan tanda $+$ atau $-$.
          Bisakah hasilnya sama dengan 0?
          \jawab{tidak bisa}{$1 + 2 + \dots + 9 = 45$ ganjil. Mengganti $+$ menjadi $-$ mengubah
          hasil sebanyak bilangan genap, jadi hasilnya selalu ganjil.}
    \item \textbf{(Buku Robek)} Sebuah buku setebal 50 lembar (100 halaman) robek satu lembarnya.
          Mungkinkah jumlah semua nomor halaman yang tersisa adalah bilangan ganjil?
          \jawab{mungkin, malah selalu ganjil}{Jumlah halaman 1 sampai 100 adalah 5050 (genap).
          Satu lembar berisi dua halaman berurutan, jumlahnya ganjil. Genap dikurangi ganjil hasilnya ganjil.}
    \item \textbf{(Koin di Meja)} Ada 9 koin yang semuanya menunjukkan sisi Angka. Setiap langkah
          kamu wajib membalik tepat 2 koin. Bisakah semua koin menunjukkan sisi Gambar?
          \jawab{tidak bisa}{Banyak koin Angka berubah 2, 0, atau $-2$, jadi tetap ganjil.
          Padahal semua Gambar artinya 0 koin Angka (genap).}
\end{enumerate}

\level{Level 4: Setara OSP (Provinsi)}
\textit{Tulis jawaban lengkap dengan alasannya.}
\begin{enumerate}
    \setcounter{enumi}{15}
    \item Di papan tertulis bilangan $1, 2, 3, \dots, 10$. Setiap langkah, kamu menghapus dua
          bilangan, lalu menuliskan selisih keduanya (bilangan besar dikurangi bilangan kecil).
          Setelah 9 langkah tersisa satu bilangan. Buktikan bahwa bilangan itu pasti ganjil.
          \jawab{bukti}{Perhatikan jumlah semua bilangan di papan. Jika $a$ dan $b$ dihapus lalu
          diganti $a - b$, jumlahnya berkurang $(a + b) - (a - b) = 2b$, yaitu bilangan genap.
          Jadi ganjil-genapnya jumlah semua bilangan \emph{tidak pernah berubah}. Di awal,
          jumlahnya $1 + 2 + \dots + 10 = 55$, ganjil. Maka di akhir, bilangan yang tersisa (yang
          sekaligus menjadi jumlahnya) juga ganjil. \hfill$\blacksquare$\\
          \emph{Ide ini disebut invarian: mencari sesuatu yang tidak berubah.}}
\end{enumerate}

\level{Level 5: Setara OSN (Nasional)}
\textit{Tulis jawaban lengkap dengan alasannya.}
\begin{enumerate}
    \setcounter{enumi}{16}
    \item Tentukan semua bilangan asli $n$ sehingga tanda $\pm$ pada
          $1 \pm 2 \pm 3 \pm \dots \pm n$ bisa dipilih ($+$ atau $-$) agar hasilnya 0.
          Buktikan bahwa tidak ada nilai $n$ lain yang memenuhi.
          \jawab{$n$ yang jika dibagi 4 bersisa 0 atau 3}{
          \emph{Syarat perlu.} Seperti soal nomor 13, hasilnya selalu sama ganjil-genapnya dengan
          $1 + 2 + \dots + n = \frac{n(n+1)}{2}$. Supaya bisa 0, nilai ini harus genap, artinya
          $n(n+1)$ habis dibagi 4. Karena $n$ dan $n+1$ tidak mungkin sama-sama genap, salah satunya
          harus habis dibagi 4. Jadi $n$ bersisa 0 atau 3 jika dibagi 4.\\
          \emph{Syarat cukup.} Empat bilangan berurutan bisa dibuat 0:
          $a - (a+1) - (a+2) + (a+3) = 0$.
          Jika $n = 4k$: kelompokkan $1, 2, \dots, n$ menjadi $k$ kelompok berisi 4, masing-masing
          dibuat 0 (kelompok pertama $1 - 2 - 3 + 4$).
          Jika $n = 4k + 3$: tulis $1 + 2 - 3 = 0$, lalu sisanya $4, 5, \dots, n$ dikelompokkan
          per 4 dengan cara yang sama. \hfill$\blacksquare$}
\end{enumerate}

\newpage
% ====================================================================
\section*{Topik 2: Pola Angka Satuan}
% ====================================================================

\begin{catatanguru}
Ayo Selidiki sebaiknya dikerjakan sendiri oleh siswa. Kesalahan yang paling sering terjadi
adalah ketika sisa pembagiannya 0 (Langkah 2).
\end{catatanguru}

\begin{selidiki}
Ambil pensilmu. Lengkapi tabel di bawah ini. Kamu cukup menulis \textbf{angka satuannya}
saja (angka paling belakang).

\begin{center}
\renewcommand{\arraystretch}{1.6}
\begin{tabular}{|l|*{8}{>{\centering\arraybackslash}p{0.9cm}|}}
\hline
Pangkat ke- & 1 & 2 & 3 & 4 & 5 & 6 & 7 & 8\\ \hline
Satuan $2^n$ & 2 & 4 & 8 & & & & & \\ \hline
Satuan $3^n$ & 3 & 9 & & & & & & \\ \hline
\end{tabular}
\end{center}

\textit{Petunjuk: kamu tidak perlu menghitung $2^8$ sampai habis. Cukup kalikan angka satuan
sebelumnya. Misalnya, satuan $2^3$ adalah 8, jadi satuan $2^4$ sama dengan satuan
$8 \times 2 = 16$, yaitu 6.}

Sekarang jawab:
\begin{enumerate}
    \item Apa yang kamu lihat pada baris $2^n$? Angka apa saja yang muncul berulang?
    \item Setelah berapa langkah polanya kembali ke awal?
    \item Tanpa menghitung, tebak angka satuan $2^{9}$ dan $2^{10}$.
\end{enumerate}

Kalau kamu menemukan bahwa angka satuannya berulang terus, selamat! Kamu baru saja menemukan
rahasia untuk menghitung bilangan raksasa seperti $3^{2026}$ tanpa kalkulator.
\end{selidiki}

\subsection*{Langkah 1: Bagi pangkatnya dengan panjang pola}
Pola angka satuan dari pangkat 3 adalah \textbf{3, 9, 7, 1}, lalu berulang. Panjang polanya 4.

\begin{contoh}
\textbf{Contoh.} Berapa angka satuan dari $3^{2026}$?

\textit{Jawab:} Bagi pangkatnya dengan 4: $2026 \div 4 = 506$ sisa \textbf{2}. Ambil angka
urutan ke-2 dari pola (3, \textbf{9}, 7, 1). Jadi angka satuannya \textbf{9}.
\end{contoh}

\begin{giliranmu}
\begin{enumerate}[label=\alph*.]
    \item Angka satuan dari $2^{15}$ adalah \kosong. \quad (pola 2: 2, 4, 8, 6)
    \item Angka satuan dari $7^{42}$ adalah \kosong. \quad (pola 7: 7, 9, 3, 1)
\end{enumerate}
\end{giliranmu}

\subsection*{Langkah 2: Hati-hati kalau sisanya 0}
Kalau pangkatnya habis dibagi panjang pola, tidak ada ``urutan ke-0''. Ambil angka yang
\textbf{paling terakhir} dari pola.

\begin{contoh}
\textbf{Contoh.} Berapa angka satuan dari $3^{100}$?

\textit{Jawab:} $100 \div 4 = 25$ sisa \textbf{0}. Ambil angka terakhir dari pola
(3, 9, 7, \textbf{1}). Jadi angka satuannya \textbf{1}.
\end{contoh}

\begin{giliranmu}
\begin{enumerate}[label=\alph*.]
    \item Angka satuan dari $8^{100}$ adalah \kosong. \quad (pola 8: 8, 4, 2, 6)
    \item Angka satuan dari $2^{2024}$ adalah \kosong.
\end{enumerate}
\end{giliranmu}

\subsection*{Langkah 3: Angka di depan tidak penting}
Angka satuan hasil perkalian hanya ditentukan oleh angka satuan bilangannya. Jadi angka
ribuan, ratusan, dan puluhan boleh kita abaikan.

\begin{contoh}
\textbf{Contoh.} Berapa angka satuan dari $2027^{2026}$?

\textit{Jawab:} Angka satuan 2027 adalah 7, jadi cukup cari angka satuan $7^{2026}$.
$2026 \div 4$ sisa 2. Urutan ke-2 dari pola (7, \textbf{9}, 3, 1) adalah \textbf{9}.
\end{contoh}

\begin{giliranmu}
\begin{enumerate}[label=\alph*.]
    \item Angka satuan dari $1234^{5}$ adalah \kosong.
    \item Angka satuan dari $2029^{2026}$ adalah \kosong.
\end{enumerate}
\end{giliranmu}

\subsection*{Langkah 4: Menjumlah dan mengali}
Cari angka satuan masing-masing dulu, baru dijumlah atau dikali. Setelah itu, ambil lagi
angka satuannya.

\begin{contoh}
\textbf{Contoh.} Berapa angka satuan dari $4^{2025} + 6^{2025}$?

\textit{Jawab:} Pola 4 adalah 4, 6, 4, 6, \dots{} Pangkatnya ganjil, jadi satuannya 4.
Pangkat dari 6 selalu berakhiran 6. Jadi $4 + 6 = 10$, dan angka satuannya \textbf{0}.
\end{contoh}

\begin{giliranmu}
\begin{enumerate}[label=\alph*.]
    \item Angka satuan dari $2^{2025} + 2^{2026}$ adalah \kosong.
    \item Angka satuan dari $2^{10} \times 3^{10}$ adalah \kosong.
\end{enumerate}
\end{giliranmu}

\begin{rangkuman}
\begin{enumerate}
    \item Tulis beberapa angka satuan awal sampai polanya berulang.
    \item Bagi pangkatnya dengan panjang pola, lalu lihat sisanya.
    \item Sisa 1 berarti urutan ke-1, sisa 2 berarti urutan ke-2, dan seterusnya.
          \textbf{Sisa 0 berarti urutan terakhir.}
    \item Angka di depan tidak penting. Cukup lihat angka satuannya.
    \item Untuk menjumlah atau mengali, cari angka satuan masing-masing dulu.
\end{enumerate}
\end{rangkuman}

\textit{Kunci Giliranmu:} Langkah 1: 8, 9. \quad Langkah 2: 6, 6. \quad Langkah 3: 4, 1. \quad
Langkah 4: 6, 6.

\newpage
\subsection*{Latihan Soal}

\level{Level 1: Pemanasan}
\begin{enumerate}
    \item Berapakah angka satuan dari $9^{2026}$?
          \jawab{1}{Pola 9, 1. Pangkat genap.}
    \item \textbf{(Bakteri Super)} Sebuah bakteri membelah diri menjadi 3 setiap jam. Pada jam
          pertama ada 3 bakteri, jam kedua 9 bakteri, dan seterusnya. Berapakah angka terakhir dari
          banyaknya bakteri pada jam ke-50?
          \jawab{9}{$3^{50}$, dan $50 \div 4$ sisa 2.}
    \item \textbf{(Mesin Tiket)} Sebuah mesin mencetak nomor seri $7^n$ untuk pembeli ke-$n$.
          Nomor seri pembeli ke-41 berakhiran angka berapa?
          \jawab{7}{$41 \div 4$ sisa 1.}
    \item Berapakah angka satuan dari $5^{2026}$ dan $6^{2026}$?
          \jawab{5 dan 6}{Pangkat 5 selalu berakhiran 5, pangkat 6 selalu berakhiran 6.}
    \item Berapakah angka satuan dari $4^{99}$?
          \jawab{4}{Pola 4, 6. Pangkat ganjil.}
\end{enumerate}

\level{Level 2: Eksplorasi}
\begin{enumerate}
    \setcounter{enumi}{5}
    \item Tentukan angka satuan dari $3^{10} + 7^{10}$.
          \jawab{8}{$9 + 9 = 18$.}
    \item \textbf{(Gembok Koper Budi)} Budi lupa kode 3 digit gembok kopernya. Petunjuknya:
          \textit{``Digit pertama adalah angka satuan $2^{40}$, digit kedua angka satuan $3^{41}$,
          dan digit ketiga angka satuan $4^{42}$.''} Berapa kodenya?
          \jawab{636}{$2^{40} \to 6$, $3^{41} \to 3$, $4^{42} \to 6$.}
    \item \textbf{(Sandi Agen 007)} PIN brankas adalah angka satuan dari \textit{``banyak hari
          dalam seminggu dipangkatkan tahun kemerdekaan Indonesia''}. Berapakah PIN-nya?
          \jawab{7}{$7^{1945}$, dan $1945 \div 4$ sisa 1.}
    \item Tentukan angka satuan dari $9^{99} - 4^{44}$.
          \jawab{3}{$9 - 6$.}
    \item \textbf{(Kalkulator Rusak)} Layar kalkulator Andi hanya menampilkan 1 digit terakhir.
          Andi mengetik $2^{2025} + 2^{2026} + 2^{2027}$. Angka berapa yang muncul?
          \jawab{4}{$2 + 4 + 8 = 14$.}
\end{enumerate}

\level{Level 3: Tantangan OSK}
\begin{enumerate}
    \setcounter{enumi}{10}
    \item Tentukan angka satuan dari $2023^{2024} + 2024^{2023}$.
          \jawab{5}{$3^{2024} \to 1$ dan $4^{2023} \to 4$.}
    \item \textbf{(Harta Karun Raja)} Raja memberi ksatrianya koin emas sebanyak
          $13^{13} \times 17^{17}$. Koin itu dibagi rata ke 10 panti asuhan. Berapa koin yang
          tersisa? \textit{(Petunjuk: sisa bagi 10 sama dengan angka satuan!)}
          \jawab{1}{$3^{13} \to 3$ dan $7^{17} \to 7$. $3 \times 7 = 21$.}
    \item Tentukan angka satuan dari $1^1 + 2^2 + 3^3 + 4^4 + 5^5$.
          \jawab{3}{$1 + 4 + 7 + 6 + 5 = 23$.}
    \item \textbf{(Bom Waktu)} Untuk menjinakkan bom, kamu harus memasukkan angka satuan dari
          $2^{3^4}$. Angka berapa yang kamu masukkan? (Kerjakan pangkat yang di atas dulu ya!)
          \jawab{2}{$3^4 = 81$, jadi $2^{81}$. $81 \div 4$ sisa 1.}
    \item Tentukan angka satuan dari $1^{2026} + 2^{2026} + 3^{2026} + \dots + 9^{2026}$.
          \jawab{5}{$1 + 4 + 9 + 6 + 5 + 6 + 9 + 4 + 1 = 45$.}
\end{enumerate}

\level{Level 4: Setara OSP (Provinsi)}
\textit{Tulis jawaban lengkap dengan alasannya.}
\begin{enumerate}
    \setcounter{enumi}{15}
    \item Buktikan bahwa untuk setiap bilangan asli $n$, bilangan $n^5 - n$ selalu habis
          dibagi 10.
          \jawab{bukti}{Cukup ditunjukkan bahwa angka satuan $n^5$ sama dengan angka satuan $n$.
          Angka satuan $n^5$ hanya bergantung pada angka satuan $n$, jadi kita cek 10 kemungkinan:
          $0^5 \to 0$, $1^5 \to 1$, $2^5 = 32 \to 2$, $3^5 = 243 \to 3$, $4^5 = 1024 \to 4$,
          $5^5 \to 5$, $6^5 \to 6$, $7^5 = 16807 \to 7$, $8^5 = 32768 \to 8$, $9^5 = 59049 \to 9$.
          Semuanya sama. Jadi $n^5$ dan $n$ punya angka satuan yang sama, sehingga $n^5 - n$
          berakhiran 0 dan habis dibagi 10. \hfill$\blacksquare$}
\end{enumerate}

\level{Level 5: Setara OSN (Nasional)}
\textit{Tulis jawaban lengkap dengan alasannya.}
\begin{enumerate}
    \setcounter{enumi}{16}
    \item Tentukan semua bilangan asli $n$ sehingga $n^n$ berakhiran angka 7. Buktikan bahwa
          tidak ada bilangan lain yang memenuhi.
          \jawab{$n$ yang jika dibagi 20 bersisa 3 atau 17}{
          Angka satuan $n^n$ ditentukan oleh angka satuan $n$. Dari pola-pola angka satuan, angka 7
          hanya bisa muncul jika angka satuan $n$ adalah 3 (pola 3, 9, 7, 1) atau 7 (pola 7, 9, 3, 1).
          Angka satuan lain tidak pernah menghasilkan 7.\\
          \emph{Kasus satuan 3:} butuh pangkat yang bersisa 3 jika dibagi 4. Tulis $n = 10k + 3$.
          Karena $10k + 3 = 8k + (2k + 3)$, sisanya jika dibagi 4 sama dengan sisa $2k + 3$. Ini
          bersisa 3 hanya jika $k$ genap, yaitu $n = 3, 23, 43, \dots$ (bersisa 3 jika dibagi 20).\\
          \emph{Kasus satuan 7:} butuh pangkat yang bersisa 1 jika dibagi 4. Tulis $n = 10k + 7$,
          sisanya sama dengan sisa $2k + 7$, yang bersisa 1 hanya jika $k$ ganjil, yaitu
          $n = 17, 37, 57, \dots$ (bersisa 17 jika dibagi 20). \hfill$\blacksquare$}
\end{enumerate}

\newpage
% ====================================================================
\section*{Topik 3: Sisa Pembagian}
% ====================================================================

\begin{catatanguru}
Hubungkan dengan Topik 2: angka satuan sama dengan sisa pembagian oleh 10. Untuk kelas 8--9 boleh
dikenalkan notasi $a \equiv b \pmod{n}$, tetapi tidak wajib.
\end{catatanguru}

\begin{selidiki}
Hari ini hari \textbf{Senin}. Lengkapi tabel berikut.

\begin{center}
\renewcommand{\arraystretch}{1.6}
\begin{tabular}{|l|*{6}{>{\centering\arraybackslash}p{1.3cm}|}}
\hline
Berapa hari lagi? & 7 & 14 & 15 & 16 & 21 & 23\\ \hline
Hari apa? & & & & & & \\ \hline
\end{tabular}
\end{center}

Sekarang jawab:
\begin{enumerate}
    \item Hari apa saja yang jatuh 7, 14, dan 21 hari lagi? Kenapa?
    \item 15 hari lagi sama dengan berapa hari lagi? Bagaimana dengan 23 hari lagi?
    \item Tanpa melihat kalender, tebak hari apa 100 hari lagi.
\end{enumerate}

Setiap 7 hari, harinya kembali sama. Jadi yang penting hanya \textbf{sisa} pembagian oleh 7.
Karena $100 = 7 \times 14 + 2$, maka 100 hari lagi sama dengan 2 hari lagi, yaitu \textbf{Rabu}.
\end{selidiki}

\subsection*{Langkah 1: Sisa pembagian}
Jika $a$ dibagi $b$, maka $a = b \times \text{(hasil bagi)} + \text{sisa}$, dan sisanya selalu
lebih kecil dari $b$. Contoh: $100 = 7 \times 14 + 2$, jadi 100 dibagi 7 bersisa 2.

\begin{contoh}
\textbf{Contoh.} Berapa sisa pembagian 2027 oleh 7?

\textit{Jawab:} Cari kelipatan 7 yang dekat: $7 \times 289 = 2023$. Jadi $2027 = 2023 + 4$,
sisanya \textbf{4}.
\end{contoh}

\begin{giliranmu}
\begin{enumerate}[label=\alph*.]
    \item Hari ini hari Kamis. 30 hari lagi adalah hari \kosong.
    \item Sisa pembagian 2026 oleh 7 adalah \kosong.
\end{enumerate}
\end{giliranmu}

\subsection*{Langkah 2: Ganti bilangan dengan sisanya}
Untuk mencari sisa dari hasil penjumlahan atau perkalian, kita boleh mengganti setiap bilangan
dengan sisanya dulu. Bilangannya boleh besar, yang dihitung hanya sisanya yang kecil.

\begin{contoh}
\textbf{Contoh.} Tentukan sisa pembagian $2025 \times 2026$ oleh 7.

\textit{Jawab:} Karena $2023 = 7 \times 289$, maka 2025 bersisa 2 dan 2026 bersisa 3.
Sisanya $2 \times 3 = \textbf{6}$.
\end{contoh}

\begin{giliranmu}
\begin{enumerate}[label=\alph*.]
    \item Sisa pembagian $23 + 19 + 18$ oleh 5 adalah \kosong.
    \item Sisa pembagian $26 \times 27$ oleh 5 adalah \kosong.
\end{enumerate}
\end{giliranmu}

\subsection*{Langkah 3: Sisa dari bilangan berpangkat}
Untuk pangkat, cari pola sisanya, sama seperti pola angka satuan di Topik 2.

\begin{contoh}
\textbf{Contoh.} Tentukan sisa pembagian $2^{100}$ oleh 7.

\textit{Jawab:} Sisa $2^1, 2^2, 2^3, 2^4, \dots$ jika dibagi 7 adalah 2, 4, 1, 2, 4, 1, \dots{}
Polanya 3 langkah. $100 \div 3$ sisa 1, jadi ambil urutan ke-1: sisanya \textbf{2}.
\end{contoh}

\begin{giliranmu}
\begin{enumerate}[label=\alph*.]
    \item Sisa pembagian $3^{50}$ oleh 4 adalah \kosong.
    \item Sisa pembagian $2^{20}$ oleh 5 adalah \kosong.
\end{enumerate}
\end{giliranmu}

\subsection*{Langkah 4: Sisa $-1$}
Kalau bilangannya kurang 1 dari pembagi (misalnya 6 dibagi 7), anggap saja sisanya $-1$.
Karena $(-1) \times (-1) = 1$, pangkatnya tinggal dilihat ganjil atau genap.

\begin{contoh}
\textbf{Contoh.} Tentukan sisa pembagian $6^{2027}$ oleh 7.

\textit{Jawab:} 6 kita anggap $-1$. Pangkatnya ganjil, jadi hasilnya $-1$. Sisa $-1$ artinya
kurang 1 dari 7, yaitu \textbf{6}.
\end{contoh}

\begin{giliranmu}
\begin{enumerate}[label=\alph*.]
    \item Sisa pembagian $4^{101}$ oleh 5 adalah \kosong.
    \item Sisa pembagian $10^{2026}$ oleh 11 adalah \kosong.
\end{enumerate}
\end{giliranmu}

\begin{rangkuman}
\begin{enumerate}
    \item Untuk soal hari, jam, atau pola berulang, yang penting hanya sisanya.
    \item Ganti setiap bilangan dengan sisanya sebelum dijumlah atau dikali.
    \item Untuk pangkat, cari pola sisanya.
    \item Kalau bilangannya kurang 1 dari pembagi, anggap sisanya $-1$.
\end{enumerate}
\end{rangkuman}

\textit{Kunci Giliranmu:} Langkah 1: Sabtu, 3. \quad Langkah 2: 0, 2. \quad Langkah 3: 1, 1.
\quad Langkah 4: 4, 1.

\newpage
\subsection*{Latihan Soal}

\level{Level 1: Pemanasan}
\begin{enumerate}
    \item Berapakah sisa pembagian 2026 oleh 9?
          \jawab{1}{$2025 = 9 \times 225$.}
    \item Hari ini hari Senin. Hari apakah 50 hari lagi?
          \jawab{Selasa}{50 dibagi 7 sisa 1.}
    \item \textbf{(Jam Dinding)} Sekarang pukul 09.00. Pukul berapakah 100 jam lagi?
          \jawab{13.00}{100 dibagi 24 sisa 4.}
    \item Tentukan sisa pembagian $25 + 38 + 47$ oleh 5 tanpa menjumlahkannya dulu.
          \jawab{0}{Sisanya $0 + 3 + 2 = 5$, habis dibagi 5.}
    \item \textbf{(Kalung Manik-Manik)} Sebuah kalung disusun dengan pola Merah, Putih, Biru, Hijau,
          Merah, Putih, Biru, Hijau, dan seterusnya. Apa warna manik ke-2026?
          \jawab{Putih}{2026 dibagi 4 sisa 2.}
\end{enumerate}

\level{Level 2: Eksplorasi}
\begin{enumerate}
    \setcounter{enumi}{5}
    \item Tentukan sisa pembagian $2024 \times 2025 \times 2026$ oleh 7.
          \jawab{6}{Sisanya $1 \times 2 \times 3$.}
    \item Tentukan sisa pembagian $2^{100}$ oleh 3.
          \jawab{1}{2 dianggap $-1$, pangkat genap.}
    \item Tentukan sisa pembagian $3^{2026}$ oleh 5.
          \jawab{4}{Angka satuan $3^{2026}$ adalah 9 (Topik 2), dan 9 dibagi 5 sisa 4.}
    \item Bilangan $n$ jika dibagi 7 bersisa 5. Tentukan sisa pembagian $n^2 + 3n$ oleh 7.
          \jawab{5}{Ganti $n$ dengan 5: $25 + 15 = 40$, dan 40 dibagi 7 sisa 5.}
    \item \textbf{(Kalender Kemerdekaan)} Tanggal 17 Agustus 2026 jatuh pada hari Senin.
          Hari apakah 17 Agustus 2028? (Ingat, 2028 adalah tahun kabisat!)
          \jawab{Kamis}{Dari 2026 ke 2027 ada 365 hari (maju 1 hari), dari 2027 ke 2028 ada 366 hari
          (maju 2 hari).}
\end{enumerate}

\level{Level 3: Tantangan OSK}
\begin{enumerate}
    \setcounter{enumi}{10}
    \item Tentukan sisa pembagian $2^{2026}$ oleh 7.
          \jawab{2}{Pola sisa 2, 4, 1. $2026 \div 3$ sisa 1.}
    \item Hari ini hari Minggu. Hari apakah $10^{10}$ hari lagi?
          \jawab{Kamis}{10 dibagi 7 sisa 3. Pola sisa $3^n$: 3, 2, 6, 4, 5, 1 (berulang setiap 6).
          $10 \div 6$ sisa 4, jadi sisanya 4. Minggu ditambah 4 hari.}
    \item Tentukan sisa pembagian $6^{2025} + 8^{2025}$ oleh 7.
          \jawab{0}{6 dianggap $-1$, 8 bersisa 1. Jadi $-1 + 1 = 0$.}
    \item Buktikan bahwa $n^3 - n$ selalu habis dibagi 3 untuk setiap bilangan bulat $n$.
          \textit{(Petunjuk: $n$ dibagi 3 sisanya hanya mungkin 0, 1, atau 2. Coba ketiganya!)}
          \jawab{bukti}{Sisa 0: $0 - 0 = 0$. Sisa 1: $1 - 1 = 0$. Sisa 2: $8 - 2 = 6$, habis dibagi 3.}
    \item Tunjukkan bahwa bilangan kuadrat jika dibagi 4 sisanya hanya 0 atau 1. Lalu, mungkinkah
          $2027 = a^2 + b^2$ untuk suatu bilangan bulat $a$ dan $b$?
          \jawab{tidak mungkin}{Bilangan genap kuadrat habis dibagi 4, bilangan ganjil kuadrat
          bersisa 1. Jadi $a^2 + b^2$ bersisa 0, 1, atau 2, padahal 2027 bersisa 3.}
\end{enumerate}

\level{Level 4: Setara OSP (Provinsi)}
\textit{Tulis jawaban lengkap dengan alasannya.}
\begin{enumerate}
    \setcounter{enumi}{15}
    \item Buktikan bahwa jika $a^2 + b^2$ habis dibagi 3, maka $a$ dan $b$ keduanya habis
          dibagi 3.
          \jawab{bukti}{Pertama, cek sisa bilangan kuadrat jika dibagi 3. Jika $a$ bersisa 0, 1,
          atau 2, maka $a^2$ bersisa $0$, $1$, atau $4 \to 1$. Jadi kuadrat hanya bersisa 0 atau 1,
          dan bersisa 0 hanya jika $a$ habis dibagi 3.\\
          Sisa $a^2 + b^2$ adalah $0 + 0$, $0 + 1$, $1 + 0$, atau $1 + 1 = 2$. Yang habis dibagi 3
          hanya $0 + 0$, artinya $a^2$ dan $b^2$ sama-sama bersisa 0. Jadi $a$ dan $b$ keduanya
          habis dibagi 3. \hfill$\blacksquare$}
\end{enumerate}

\level{Level 5: Setara OSN (Nasional)}
\textit{Tulis jawaban lengkap dengan alasannya.}
\begin{enumerate}
    \setcounter{enumi}{16}
    \item Tentukan semua bilangan asli $n$ sehingga $n \times 2^n + 1$ habis dibagi 3.
          Buktikan jawabanmu.
          \jawab{$n$ yang jika dibagi 6 bersisa 1 atau 2}{
          Jika dibagi 3, angka 2 dianggap $-1$. Jadi $2^n$ bersisa 1 jika $n$ genap dan bersisa
          $-1$ jika $n$ ganjil.\\
          \emph{Kasus $n$ genap:} $n \times 2^n + 1$ bersisa sama dengan $n + 1$. Habis dibagi 3
          jika $n$ bersisa 2 jika dibagi 3. Bilangan genap yang bersisa 2 jika dibagi 3 adalah
          2, 8, 14, \dots, yaitu yang bersisa 2 jika dibagi 6.\\
          \emph{Kasus $n$ ganjil:} bersisa sama dengan $-n + 1$. Habis dibagi 3 jika $n$ bersisa 1
          jika dibagi 3. Bilangan ganjil seperti itu adalah 1, 7, 13, \dots, yaitu yang bersisa 1
          jika dibagi 6.\\
          Jadi jawabannya semua $n$ yang bersisa 1 atau 2 jika dibagi 6. \hfill$\blacksquare$}
\end{enumerate}

\newpage
% ====================================================================
\section*{Topik 4: Ciri Habis Dibagi}
% ====================================================================

\begin{catatanguru}
Setelah Ayo Selidiki, tunjukkan alasan ciri habis dibagi 9: $10, 100, 1000, \dots$ semuanya
bersisa 1 jika dibagi 9 (Topik 3). Tekankan jebakan pembagi gabungan: 24 bukan $4 \times 6$.
\end{catatanguru}

\begin{selidiki}
Semua bilangan di bawah ini adalah kelipatan 9. Jumlahkan digit-digitnya.

\begin{center}
\renewcommand{\arraystretch}{1.6}
\begin{tabular}{|l|*{7}{>{\centering\arraybackslash}p{1.1cm}|}}
\hline
Bilangan & 18 & 27 & 54 & 99 & 117 & 315 & 2025\\ \hline
Jumlah digit & 9 & & & & & & \\ \hline
\end{tabular}
\end{center}

Sekarang jawab:
\begin{enumerate}
    \item Apa yang kamu lihat pada jumlah digitnya?
    \item Tanpa membagi, tebak: apakah 1233 habis dibagi 9? Bagaimana dengan 1234?
    \item Coba juga untuk kelipatan 3: 12, 15, 21, 111. Apa yang kamu lihat?
\end{enumerate}

Jumlah digit kelipatan 9 selalu habis dibagi 9. Ini adalah \textbf{ciri habis dibagi 9}.
Masih ada banyak ciri lain yang akan kita pelajari.
\end{selidiki}

\subsection*{Langkah 1: Ciri dari digit terakhir}
Sebuah bilangan habis dibagi \textbf{2} jika digit terakhirnya genap, habis dibagi \textbf{5}
jika digit terakhirnya 0 atau 5, habis dibagi \textbf{4} jika 2 digit terakhirnya habis dibagi 4,
dan habis dibagi \textbf{8} jika 3 digit terakhirnya habis dibagi 8.

\begin{contoh}
\textbf{Contoh.} Apakah 123456 habis dibagi 4? Apakah habis dibagi 8?

\textit{Jawab:} Dua digit terakhir 56, dan $56 = 4 \times 14$, jadi habis dibagi 4.
Tiga digit terakhir 456, dan $456 = 8 \times 57$, jadi habis dibagi 8 juga.
\end{contoh}

\begin{giliranmu}
\begin{enumerate}[label=\alph*.]
    \item Apakah 2026 habis dibagi 4? \kosong
    \item Apakah 51128 habis dibagi 8? \kosong
\end{enumerate}
\end{giliranmu}

\subsection*{Langkah 2: Ciri dari jumlah digit}
Sebuah bilangan habis dibagi \textbf{3} jika jumlah digitnya habis dibagi 3, dan habis dibagi
\textbf{9} jika jumlah digitnya habis dibagi 9. Bonus: sisa pembagian oleh 9 sama dengan sisa
jumlah digitnya dibagi 9.

\begin{contoh}
\textbf{Contoh.} Apakah 123456789 habis dibagi 9?

\textit{Jawab:} Jumlah digitnya $1 + 2 + \dots + 9 = 45$, dan 45 habis dibagi 9. Jadi \textbf{ya}.
\end{contoh}

\begin{giliranmu}
\begin{enumerate}[label=\alph*.]
    \item Apakah 2025 habis dibagi 9? \kosong
    \item Sisa pembagian 1234 oleh 9 adalah \kosong.
\end{enumerate}
\end{giliranmu}

\subsection*{Langkah 3: Ciri habis dibagi 11}
Jumlahkan digit secara selang-seling dari kanan: tambah, kurang, tambah, kurang, dan seterusnya.
Jika hasilnya habis dibagi 11 (termasuk 0), bilangan itu habis dibagi 11.

\begin{contoh}
\textbf{Contoh.} Apakah 2728 habis dibagi 11?

\textit{Jawab:} Dari kanan: $8 - 2 + 7 - 2 = 11$. Karena 11 habis dibagi 11, maka 2728 habis
dibagi 11. (Memang $2728 = 11 \times 248$.)
\end{contoh}

\begin{giliranmu}
\begin{enumerate}[label=\alph*.]
    \item Apakah 1331 habis dibagi 11? \kosong
    \item Apakah 2026 habis dibagi 11? \kosong
\end{enumerate}
\end{giliranmu}

\subsection*{Langkah 4: Pembagi gabungan}
Habis dibagi 6 artinya habis dibagi 2 \textbf{dan} 3. Habis dibagi 72 artinya habis dibagi 8
\textbf{dan} 9. Hati-hati, kedua pembaginya tidak boleh punya faktor yang sama. Contohnya, 12
habis dibagi 4 dan 6, tetapi tidak habis dibagi 24.

\begin{contoh}
\textbf{Contoh.} Bilangan 5 digit $7a36b$ habis dibagi 72. Tentukan semua kemungkinannya.

\textit{Jawab:} Supaya habis dibagi 8, tiga digit terakhir ($36b$) harus habis dibagi 8, jadi
$b = 0$ atau $b = 8$. Supaya habis dibagi 9, jumlah digit $16 + a + b$ harus kelipatan 9.
Kalau $b = 0$, maka $a = 2$. Kalau $b = 8$, maka $a = 3$. Jadi bilangannya \textbf{72360}
atau \textbf{73368}.
\end{contoh}

\begin{giliranmu}
\begin{enumerate}[label=\alph*.]
    \item Apakah 2028 habis dibagi 6? \kosong
    \item Digit $a$ supaya $3a6$ habis dibagi 6 adalah \kosong.
\end{enumerate}
\end{giliranmu}

\begin{rangkuman}
\begin{enumerate}
    \item Dibagi 2, 5: lihat 1 digit terakhir. Dibagi 4: 2 digit terakhir. Dibagi 8: 3 digit terakhir.
    \item Dibagi 3 dan 9: lihat jumlah semua digit.
    \item Dibagi 11: jumlahkan digit selang-seling (tambah, kurang, tambah, \dots).
    \item Pembagi gabungan: pecah menjadi dua pembagi yang tidak punya faktor sama.
\end{enumerate}
\end{rangkuman}

\textit{Kunci Giliranmu:} Langkah 1: tidak, ya. \quad Langkah 2: ya, 1. \quad Langkah 3: ya, tidak.
\quad Langkah 4: ya, 0, 3, 6, atau 9.

\newpage
\subsection*{Latihan Soal}

\level{Level 1: Pemanasan}
\begin{enumerate}
    \item Manakah yang habis dibagi 3: 2025, 1234, 111111, 2026?
          \jawab{2025 dan 111111}{Jumlah digitnya 9 dan 6.}
    \item Apakah 7392 habis dibagi 8?
          \jawab{ya}{$392 = 8 \times 49$.}
    \item Apakah 918082 habis dibagi 11?
          \jawab{ya}{$2 - 8 + 0 - 8 + 1 - 9 = -22$.}
    \item Bilangan $45\square$ habis dibagi 6. Angka berapa saja yang bisa mengisi kotak?
          \jawab{0 atau 6}{Harus genap dan $4 + 5 + \square$ habis dibagi 3.}
    \item Tentukan sisa pembagian 20262026 oleh 9.
          \jawab{2}{Jumlah digitnya 20, dan 20 dibagi 9 sisa 2.}
\end{enumerate}

\level{Level 2: Eksplorasi}
\begin{enumerate}
    \setcounter{enumi}{5}
    \item Tentukan digit $a$ supaya bilangan $2a26$ habis dibagi 11.
          \jawab{$a = 9$}{$6 - 2 + a - 2 = a + 2$ harus habis dibagi 11.}
    \item Ada berapa bilangan 4 digit berbentuk $1a2b$ yang habis dibagi 15?
          \jawab{7}{$b = 0$: $a$ = 0, 3, 6, 9. $b = 5$: $a$ = 1, 4, 7.}
    \item \textbf{(Nota yang Robek)} Pak Ali membeli 72 pensil yang harganya sama. Total harganya
          tertulis Rp$\square679\square$, tetapi digit pertama dan terakhirnya terhapus air hujan.
          Berapakah harga satu pensil?
          \jawab{Rp511}{Habis dibagi 8: $79\square$ harus kelipatan 8, jadi digit terakhir 2.
          Habis dibagi 9: digit pertama 3. Total Rp36.792, dibagi 72.}
    \item Buktikan bahwa selisih bilangan dua digit dengan bilangan kebalikannya (misalnya
          $73 - 37$) selalu habis dibagi 9.
          \jawab{bukti}{Bilangan $ab$ bernilai $10a + b$. Selisihnya $(10a + b) - (10b + a) = 9a - 9b$.}
    \item Berapakah bilangan terkecil yang semua digitnya angka 1 dan habis dibagi 9?
          \jawab{111111111}{Butuh 9 angka 1 supaya jumlah digitnya 9.}
\end{enumerate}

\level{Level 3: Tantangan OSK}
\begin{enumerate}
    \setcounter{enumi}{10}
    \item Berapakah bilangan terkecil yang semua digitnya angka 1 dan habis dibagi 99?
          \jawab{18 angka 1}{Harus habis dibagi 9 (banyak angka 1 kelipatan 9) dan habis dibagi 11
          (banyak angka 1 harus genap).}
    \item Bilangan prima terbesar berapakah yang pasti membagi semua bilangan 6 digit berbentuk
          $abcabc$?
          \jawab{13}{$abcabc = abc \times 1001 = abc \times 7 \times 11 \times 13$.}
    \item Ada berapa bilangan 3 digit yang habis dibagi 3 tetapi tidak habis dibagi 5?
          \jawab{240}{Kelipatan 3 ada 300, kelipatan 15 ada 60.}
    \item Bilangan $N$ dibentuk dengan menulis ``2026'' sebanyak 100 kali berturut-turut.
          Tentukan sisa pembagian $N$ oleh 9.
          \jawab{1}{Jumlah digitnya $10 \times 100 = 1000$, dan 1000 dibagi 9 sisa 1.}
    \item Tentukan semua bilangan 6 digit $a2026b$ yang habis dibagi 36.
          \jawab{420264, 820260, 920268}{Habis dibagi 4: $6b$ kelipatan 4, jadi $b$ = 0, 4, atau 8.
          Habis dibagi 9: $a + b + 10$ kelipatan 9.}
\end{enumerate}

\level{Level 4: Setara OSP (Provinsi)}
\textit{Tulis jawaban lengkap dengan alasannya.}
\begin{enumerate}
    \setcounter{enumi}{15}
    \item Buktikan bahwa bilangan tiga digit $abc$ habis dibagi 9 jika dan hanya jika
          $a + b + c$ habis dibagi 9.
          \jawab{bukti}{Bilangan $abc$ bernilai $100a + 10b + c$. Tulis
          \[ 100a + 10b + c = (99a + 9b) + (a + b + c) = 9(11a + b) + (a + b + c). \]
          Bagian $9(11a + b)$ selalu habis dibagi 9. Jadi $abc$ dan $a + b + c$ selalu punya sisa
          yang sama jika dibagi 9. Akibatnya, $abc$ habis dibagi 9 tepat ketika $a + b + c$ habis
          dibagi 9. \hfill$\blacksquare$}
\end{enumerate}

\level{Level 5: Setara OSN (Nasional)}
\textit{Tulis jawaban lengkap dengan alasannya.}
\begin{enumerate}
    \setcounter{enumi}{16}
    \item Tentukan semua bilangan 4 digit berbentuk $aabb$ (dua digit pertama sama dan dua digit
          terakhir sama) yang merupakan bilangan kuadrat. Buktikan tidak ada yang lain.
          \jawab{7744}{$aabb = 1100a + 11b = 11(100a + b)$. Karena bilangan kuadrat dan habis dibagi
          11 (prima), ia harus habis dibagi $11^2$, jadi $100a + b$ habis dibagi 11. Dengan ciri
          habis dibagi 11, bilangan $a0b$ memberi $b - 0 + a = a + b$, sehingga $a + b$ habis dibagi
          11. Karena $1 \le a + b \le 18$, maka $a + b = 11$.\\
          Maka $100a + b = 100a + 11 - a = 99a + 11 = 11(9a + 1)$, dan $aabb = 121(9a + 1)$. Supaya
          kuadrat, $9a + 1$ harus kuadrat. Untuk $a = 2, \dots, 9$ (karena $b = 11 - a \le 9$),
          nilai $9a + 1$ adalah 19, 28, 37, 46, 55, 64, 73, 82. Yang kuadrat hanya $64$ ($a = 7$).
          Jadi $b = 4$, dan bilangannya $7744 = 88^2$. \hfill$\blacksquare$}
\end{enumerate}

\newpage
% ====================================================================
\section*{Topik 5: Bilangan Prima dan Faktor}
% ====================================================================

\begin{catatanguru}
Contoh di Langkah 3 (loker) bisa diperagakan dengan 10 kertas bernomor di papan. Hubungkan
dengan Ayo Selidiki: bilangan yang banyak faktornya ganjil adalah bilangan kuadrat.
\end{catatanguru}

\begin{selidiki}
Tuliskan semua faktor dari setiap bilangan, lalu hitung banyaknya.

\begin{center}
\renewcommand{\arraystretch}{1.5}
\begin{tabular}{|c|p{5cm}|c|}
\hline
Bilangan & Faktor-faktornya & Banyak faktor\\ \hline
2 & 1, 2 & 2\\ \hline
4 & 1, 2, 4 & 3\\ \hline
6 & & \\ \hline
7 & & \\ \hline
9 & & \\ \hline
12 & & \\ \hline
13 & & \\ \hline
16 & & \\ \hline
\end{tabular}
\end{center}

Sekarang jawab:
\begin{enumerate}
    \item Bilangan mana saja yang faktornya tepat 2?
    \item Bilangan mana saja yang banyak faktornya ganjil? Apa istimewanya bilangan-bilangan itu?
\end{enumerate}

Bilangan yang faktornya tepat 2 (yaitu 1 dan dirinya sendiri) disebut \textbf{bilangan prima}.
Bilangan yang banyak faktornya ganjil ternyata selalu bilangan kuadrat. Kita akan lihat
kenapa di Langkah 3.
\end{selidiki}

\subsection*{Langkah 1: Bilangan prima}
Bilangan prima: 2, 3, 5, 7, 11, 13, 17, 19, 23, 29, \dots{} Angka 1 bukan prima, dan 2 adalah
satu-satunya prima yang genap. Untuk mengecek apakah sebuah bilangan prima, cukup coba bagi
dengan bilangan prima yang kuadratnya tidak lebih dari bilangan itu.

\begin{contoh}
\textbf{Contoh.} Apakah 97 bilangan prima?

\textit{Jawab:} $10^2 = 100 > 97$, jadi cukup coba bagi dengan 2, 3, 5, 7. Tidak ada yang habis.
Jadi 97 \textbf{prima}.
\end{contoh}

\begin{giliranmu}
\begin{enumerate}[label=\alph*.]
    \item Apakah 91 bilangan prima? \kosong
    \item Apakah 101 bilangan prima? \kosong
\end{enumerate}
\end{giliranmu}

\subsection*{Langkah 2: Faktorisasi prima dan banyak faktor}
Setiap bilangan bisa ditulis sebagai perkalian bilangan prima, misalnya $72 = 2^3 \times 3^2$.
Untuk menghitung banyak faktor: \textbf{tambah 1 setiap pangkatnya, lalu kalikan}.

\begin{contoh}
\textbf{Contoh.} Berapa banyak faktor dari 72?

\textit{Jawab:} $72 = 2^3 \times 3^2$. Banyak faktornya $(3 + 1) \times (2 + 1) = \textbf{12}$.
(Cek: 1, 2, 3, 4, 6, 8, 9, 12, 18, 24, 36, 72.)
\end{contoh}

\begin{giliranmu}
\begin{enumerate}[label=\alph*.]
    \item Banyak faktor dari $36$ adalah \kosong.
    \item Banyak faktor dari $48$ adalah \kosong.
\end{enumerate}
\end{giliranmu}

\subsection*{Langkah 3: Bilangan kuadrat}
Faktor biasanya berpasangan, misalnya $12 = 1 \times 12 = 2 \times 6 = 3 \times 4$. Pada bilangan
kuadrat seperti $36 = 6 \times 6$, ada faktor yang berpasangan dengan dirinya sendiri, jadi banyak
faktornya ganjil. Pada faktorisasi prima, bilangan kuadrat \textbf{semua pangkatnya genap}.

\begin{contoh}
\textbf{Contoh (Loker).} Ada 100 loker tertutup. Siswa ke-1 membalik (buka/tutup) semua loker,
siswa ke-2 membalik loker nomor 2, 4, 6, \dots, siswa ke-3 membalik loker 3, 6, 9, \dots, sampai
siswa ke-100. Berapa loker yang terbuka di akhir?

\textit{Jawab:} Loker nomor $k$ dibalik sebanyak banyak faktor $k$. Loker terbuka jika dibalik
ganjil kali, yaitu jika $k$ bilangan kuadrat: 1, 4, 9, \dots, 100. Jadi ada \textbf{10 loker}.
\end{contoh}

\begin{giliranmu}
\begin{enumerate}[label=\alph*.]
    \item Bilangan terkecil yang jika dikalikan 12 menjadi bilangan kuadrat adalah \kosong.
    \item Apakah $2^4 \times 3^2 \times 5^2$ bilangan kuadrat? \kosong
\end{enumerate}
\end{giliranmu}

\subsection*{Langkah 4: Angka nol di belakang $n!$}
Notasi $n!$ artinya $1 \times 2 \times 3 \times \dots \times n$. Setiap angka nol di belakang
berasal dari $10 = 2 \times 5$. Faktor 2 selalu lebih banyak, jadi cukup \textbf{hitung faktor 5}.

\begin{contoh}
\textbf{Contoh.} Ada berapa angka nol di bagian paling belakang $100!$?

\textit{Jawab:} Kelipatan 5 sampai 100 ada 20. Kelipatan 25 (25, 50, 75, 100) menyumbang satu
faktor 5 lagi, ada 4. Jadi ada $20 + 4 = \textbf{24}$ angka nol.
\end{contoh}

\begin{giliranmu}
\begin{enumerate}[label=\alph*.]
    \item Banyak angka nol di belakang $10!$ adalah \kosong.
    \item Banyak angka nol di belakang $50!$ adalah \kosong.
\end{enumerate}
\end{giliranmu}

\begin{rangkuman}
\begin{enumerate}
    \item Bilangan prima hanya punya 2 faktor. Angka 1 bukan prima, dan 2 satu-satunya prima genap.
    \item Banyak faktor: tambah 1 setiap pangkat pada faktorisasi prima, lalu kalikan.
    \item Bilangan kuadrat: semua pangkatnya genap, dan banyak faktornya ganjil.
    \item Angka nol di belakang $n!$: hitung banyak faktor 5.
\end{enumerate}
\end{rangkuman}

\textit{Kunci Giliranmu:} Langkah 1: tidak ($7 \times 13$), ya. \quad Langkah 2: 9, 10. \quad
Langkah 3: 3, ya. \quad Langkah 4: 2, 12.

\newpage
\subsection*{Latihan Soal}

\level{Level 1: Pemanasan}
\begin{enumerate}
    \item Tuliskan semua bilangan prima di antara 50 dan 70.
          \jawab{53, 59, 61, 67}{}
    \item Manakah yang merupakan bilangan prima: 89, 91, 97, 119?
          \jawab{89 dan 97}{$91 = 7 \times 13$ dan $119 = 7 \times 17$.}
    \item Tuliskan faktorisasi prima dari 2024 dan 2025.
          \jawab{$2^3 \times 11 \times 23$ dan $3^4 \times 5^2$}{}
    \item Berapakah banyak faktor dari 90?
          \jawab{12}{$90 = 2 \times 3^2 \times 5$, jadi $2 \times 3 \times 2$.}
    \item Ada berapa angka nol di bagian paling belakang $25!$?
          \jawab{6}{Kelipatan 5 ada 5, kelipatan 25 ada 1.}
\end{enumerate}

\level{Level 2: Eksplorasi}
\begin{enumerate}
    \setcounter{enumi}{5}
    \item Berapakah bilangan terkecil yang mempunyai tepat 6 faktor?
          \jawab{12}{$12 = 2^2 \times 3$, banyak faktornya $3 \times 2 = 6$.}
    \item Ada berapa faktor dari 360 yang merupakan bilangan genap?
          \jawab{18}{$360 = 2^3 \times 3^2 \times 5$. Pangkat 2 harus 1, 2, atau 3 (3 pilihan),
          pangkat 3 ada 3 pilihan, pangkat 5 ada 2 pilihan.}
    \item \textbf{(Kotak Cokelat)} Bu Rina punya 60 cokelat dan ingin memasukkannya ke beberapa
          kotak dengan isi yang sama. Setiap kotak berisi minimal 2 cokelat dan kotaknya minimal 2.
          Ada berapa cara?
          \jawab{10}{60 punya 12 faktor. Kurangi faktor 1 dan 60.}
    \item Berapakah bilangan terkecil yang jika dikalikan dengan 180 menjadi bilangan kuadrat?
          \jawab{5}{$180 = 2^2 \times 3^2 \times 5$.}
    \item Ada berapa angka nol di bagian paling belakang $2026!$?
          \jawab{505}{$405 + 81 + 16 + 3$ (kelipatan 5, 25, 125, dan 625).}
\end{enumerate}

\level{Level 3: Tantangan OSK}
\begin{enumerate}
    \setcounter{enumi}{10}
    \item Ada berapa bilangan dari 1 sampai 100 yang mempunyai tepat 3 faktor?
          \jawab{4}{Hanya kuadrat bilangan prima: 4, 9, 25, 49.}
    \item \textbf{(Lampu Aula)} Di aula ada 2026 lampu bernomor 1 sampai 2026, semuanya mati.
          Siswa ke-$k$ menekan saklar semua lampu yang nomornya kelipatan $k$, untuk
          $k = 1, 2, \dots, 2026$. Ada berapa lampu yang menyala di akhir?
          \jawab{45}{Sama seperti soal loker. $45^2 = 2025$.}
    \item Tentukan semua bilangan $n$ sehingga $n!$ diakhiri tepat 6 angka nol.
          \jawab{$n$ = 25, 26, 27, 28, 29}{$24!$ punya 4 nol, $25!$ punya 6, dan $30!$ punya 7.}
    \item Berapakah bilangan terkecil yang mempunyai tepat 12 faktor?
          \jawab{60}{$60 = 2^2 \times 3 \times 5$, banyak faktornya $3 \times 2 \times 2$.}
    \item Tentukan faktor prima terbesar dari $2^{16} - 1$.
          \textit{(Petunjuk: $a^2 - b^2 = (a-b)(a+b)$.)}
          \jawab{257}{$2^{16} - 1 = (2^8 - 1)(2^8 + 1) = 255 \times 257 = 3 \times 5 \times 17 \times 257$.}
\end{enumerate}

\level{Level 4: Setara OSP (Provinsi)}
\textit{Tulis jawaban lengkap dengan alasannya.}
\begin{enumerate}
    \setcounter{enumi}{15}
    \item Buktikan bahwa jika $p$ bilangan prima yang lebih dari 3, maka $p^2 - 1$ habis
          dibagi 24.
          \jawab{bukti}{Tulis $p^2 - 1 = (p - 1)(p + 1)$. Karena $24 = 3 \times 8$, cukup buktikan
          habis dibagi 3 dan habis dibagi 8.\\
          \emph{Habis dibagi 3:} di antara tiga bilangan berurutan $p - 1$, $p$, $p + 1$, salah
          satunya habis dibagi 3. Karena $p$ prima lebih dari 3, $p$ tidak habis dibagi 3. Jadi
          $p - 1$ atau $p + 1$ habis dibagi 3.\\
          \emph{Habis dibagi 8:} $p$ ganjil, jadi $p - 1$ dan $p + 1$ adalah dua bilangan genap
          berurutan. Salah satunya pasti kelipatan 4, dan yang lain kelipatan 2. Hasil kalinya
          habis dibagi $4 \times 2 = 8$. \hfill$\blacksquare$}
\end{enumerate}

\level{Level 5: Setara OSN (Nasional)}
\textit{Tulis jawaban lengkap dengan alasannya.}
\begin{enumerate}
    \setcounter{enumi}{16}
    \item Tentukan semua bilangan prima $p$ sehingga $p + 10$ dan $p + 14$ juga bilangan prima.
          Buktikan tidak ada yang lain.
          \jawab{$p = 3$}{Untuk $p = 3$: $13$ dan $17$ prima, jadi memenuhi.\\
          Sekarang lihat sisa pembagian oleh 3. Jika $p$ bersisa 1, maka $p + 14$ bersisa
          $1 + 2 = 3 \to 0$, jadi habis dibagi 3 dan lebih dari 3, bukan prima. Jika $p$ bersisa 2,
          maka $p + 10$ bersisa $2 + 1 = 3 \to 0$, juga bukan prima. Jadi $p$ harus habis dibagi 3,
          dan satu-satunya prima yang habis dibagi 3 adalah $p = 3$. \hfill$\blacksquare$}
\end{enumerate}

\newpage
% ====================================================================
\section*{Topik 6: FPB dan KPK}
% ====================================================================

\begin{catatanguru}
Kebanyakan siswa sudah kenal FPB dan KPK dari sekolah. Fokuskan pada membedakan soal cerita FPB
dan KPK, rumus FPB $\times$ KPK $= a \times b$, dan cara mencari FPB bilangan besar (Langkah 4).
\end{catatanguru}

\begin{selidiki}
Lengkapi daftar kelipatan berikut, lalu \textbf{lingkari} bilangan yang muncul di kedua baris.

\begin{center}
\renewcommand{\arraystretch}{1.6}
\begin{tabular}{|l|*{9}{>{\centering\arraybackslash}p{0.8cm}|}}
\hline
Kelipatan 4 & 4 & 8 & & & & & & & 36\\ \hline
Kelipatan 6 & 6 & 12 & & & & & & & 54\\ \hline
\end{tabular}
\end{center}

Sekarang jawab:
\begin{enumerate}
    \item Bilangan berapa saja yang kamu lingkari? Mana yang paling kecil?
    \item Tulis semua faktor 12 dan semua faktor 18. Faktor mana yang sama? Mana yang terbesar?
\end{enumerate}

Bilangan terkecil yang muncul di kedua baris disebut \textbf{KPK} (Kelipatan Persekutuan
terKecil). Faktor terbesar yang sama disebut \textbf{FPB} (Faktor Persekutuan terBesar).
\end{selidiki}

\subsection*{Langkah 1: FPB dan KPK dengan faktorisasi prima}
\textbf{FPB:} ambil faktor prima yang sama, dengan pangkat terkecil.
\textbf{KPK:} ambil semua faktor prima, dengan pangkat terbesar.

\begin{contoh}
\textbf{Contoh.} Tentukan FPB dan KPK dari 72 dan 120.

\textit{Jawab:} $72 = 2^3 \times 3^2$ dan $120 = 2^3 \times 3 \times 5$.
FPB $= 2^3 \times 3 = \textbf{24}$. KPK $= 2^3 \times 3^2 \times 5 = \textbf{360}$.
\end{contoh}

\begin{giliranmu}
\begin{enumerate}[label=\alph*.]
    \item FPB dan KPK dari 20 dan 30 adalah \kosong{} dan \kosong.
    \item FPB dan KPK dari 24 dan 36 adalah \kosong{} dan \kosong.
\end{enumerate}
\end{giliranmu}

\subsection*{Langkah 2: Soal cerita, FPB atau KPK?}
Kata ``dibagi rata sebanyak-banyaknya'' atau ``sepanjang-panjangnya'' biasanya \textbf{FPB}.
Kata ``bersamaan lagi'' atau ``bertemu lagi'' biasanya \textbf{KPK}.

\begin{contoh}
\textbf{Contoh.} Tiga lampu berkedip setiap 4, 6, dan 10 detik. Saat ini ketiganya berkedip
bersamaan. Berapa detik lagi mereka berkedip bersamaan lagi?

\textit{Jawab:} Kata kuncinya ``bersamaan lagi'', jadi KPK. KPK$(4, 6, 10) = \textbf{60}$ detik.
\end{contoh}

\begin{giliranmu}
\begin{enumerate}[label=\alph*.]
    \item 24 apel dan 36 jeruk dibagi rata ke kantong sebanyak-banyaknya. Banyak kantongnya \kosong.
    \item Ani les setiap 3 hari dan Budi setiap 4 hari. Hari ini mereka les bersama. Mereka les
          bersama lagi \kosong{} hari kemudian.
\end{enumerate}
\end{giliranmu}

\subsection*{Langkah 3: FPB $\times$ KPK $= a \times b$}
Untuk dua bilangan $a$ dan $b$, hasil kali FPB dan KPK-nya selalu sama dengan $a \times b$.

\begin{contoh}
\textbf{Contoh.} FPB dua bilangan adalah 4 dan KPK-nya 48. Salah satu bilangannya 12.
Berapa bilangan yang lain?

\textit{Jawab:} $4 \times 48 = 12 \times b$, jadi $b = 192 \div 12 = \textbf{16}$.
\end{contoh}

\begin{giliranmu}
\begin{enumerate}[label=\alph*.]
    \item FPB 5, KPK 60, dan salah satu bilangannya 15. Bilangan yang lain \kosong.
    \item Hasil kali dua bilangan 180 dan FPB-nya 6. KPK-nya \kosong.
\end{enumerate}
\end{giliranmu}

\subsection*{Langkah 4: FPB bilangan besar}
Kalau bilangannya besar dan sulit difaktorkan, gunakan cara ini:
FPB$(a, b)$ = FPB$(b,$ sisa $a \div b)$. Ulangi sampai sisanya 0.

\begin{contoh}
\textbf{Contoh.} Tentukan FPB$(1309, 1001)$.

\textit{Jawab:} 1309 dibagi 1001 sisa 308. 1001 dibagi 308 sisa 77. 308 dibagi 77 sisa 0.
Jadi FPB-nya \textbf{77}.
\end{contoh}

\begin{giliranmu}
\begin{enumerate}[label=\alph*.]
    \item FPB$(391, 299)$ adalah \kosong.
    \item FPB$(2026, 2024)$ adalah \kosong.
\end{enumerate}
\end{giliranmu}

\begin{rangkuman}
\begin{enumerate}
    \item FPB: faktor prima yang sama, pangkat terkecil. KPK: semua faktor prima, pangkat terbesar.
    \item ``Dibagi rata sebanyak-banyaknya'' biasanya FPB. ``Bersamaan lagi'' biasanya KPK.
    \item FPB $\times$ KPK $= a \times b$.
    \item FPB bilangan besar: FPB$(a, b)$ = FPB$(b,$ sisa $a \div b)$.
\end{enumerate}
\end{rangkuman}

\textit{Kunci Giliranmu:} Langkah 1: 10 dan 60, 12 dan 72. \quad Langkah 2: 12, 12. \quad
Langkah 3: 20, 30. \quad Langkah 4: 23, 2.

\newpage
\subsection*{Latihan Soal}

\level{Level 1: Pemanasan}
\begin{enumerate}
    \item Tentukan FPB dan KPK dari 12 dan 18.
          \jawab{FPB 6, KPK 36}{}
    \item Tentukan FPB dan KPK dari $2^3 \times 3^2 \times 5$ dan $2^2 \times 3^4 \times 7$.
          \jawab{FPB $2^2 \times 3^2 = 36$, KPK $2^3 \times 3^4 \times 5 \times 7 = 22680$}{}
    \item \textbf{(Bus Sekolah)} Bus A berangkat setiap 15 menit dan bus B setiap 20 menit.
          Keduanya berangkat bersama pukul 06.00. Pukul berapa keduanya berangkat bersama lagi?
          \jawab{07.00}{KPK$(15, 20) = 60$ menit.}
    \item FPB dua bilangan adalah 6 dan KPK-nya 72. Jika salah satu bilangannya 18, berapakah
          bilangan yang lain?
          \jawab{24}{$6 \times 72 \div 18$.}
    \item \textbf{(Ubin Kelas)} Lantai kelas berukuran 420 cm $\times$ 315 cm akan ditutup ubin
          persegi yang sama besar, sebesar-besarnya, tanpa dipotong. Berapa ukuran ubinnya dan
          berapa banyak ubin yang dibutuhkan?
          \jawab{ubin 105 cm, 12 ubin}{FPB$(420, 315) = 105$. Banyak ubin $4 \times 3$.}
\end{enumerate}

\level{Level 2: Eksplorasi}
\begin{enumerate}
    \setcounter{enumi}{5}
    \item \textbf{(Tiga Lonceng)} Tiga lonceng berdentang setiap 8, 12, dan 18 menit. Ketiganya
          berdentang bersama pukul 07.00. Berapa kali ketiganya berdentang bersama dari pukul 07.00
          sampai 17.00 (termasuk pukul 07.00)?
          \jawab{9 kali}{KPK = 72 menit. Dalam 600 menit: menit ke-0, 72, 144, \dots, 576.}
    \item Berapakah FPB$(2025, 2026)$? Jelaskan alasannya.
          \jawab{1}{Bilangan yang membagi keduanya juga membagi selisihnya, yaitu 1.}
    \item Carilah FPB$(1517, 1147)$.
          \jawab{37}{1517 dibagi 1147 sisa 370, 1147 dibagi 370 sisa 37, 370 dibagi 37 sisa 0.}
    \item Tentukan bilangan terkecil yang lebih dari 2 dan jika dibagi 3, 4, maupun 5 selalu
          bersisa 2.
          \jawab{62}{KPK$(3, 4, 5) = 60$, ditambah 2.}
    \item \textbf{(Roda Gigi)} Roda gigi A punya 24 gigi dan roda gigi B punya 36 gigi. Keduanya
          saling bersinggungan dan diberi tanda di titik pertemuannya. Setelah berapa putaran
          roda A kedua tanda bertemu lagi? Berapa putaran roda B saat itu?
          \jawab{A 3 putaran, B 2 putaran}{KPK$(24, 36) = 72$ gigi.}
\end{enumerate}

\level{Level 3: Tantangan OSK}
\begin{enumerate}
    \setcounter{enumi}{10}
    \item Tentukan bilangan terkecil yang jika dibagi 3 bersisa 2, dibagi 4 bersisa 3, dan dibagi
          5 bersisa 4. \textit{(Petunjuk: apa yang terjadi kalau bilangan itu ditambah 1?)}
          \jawab{59}{Bilangan itu ditambah 1 habis dibagi 3, 4, dan 5. Jadi $60 - 1$.}
    \item \textbf{(Barisan Upacara)} Peserta upacara jika dibariskan 4, 5, atau 6 orang per baris
          selalu tersisa 1 orang. Tetapi jika dibariskan 7 orang per baris, tidak ada yang tersisa.
          Berapa banyak peserta paling sedikit?
          \jawab{301}{Bilangannya berbentuk $60k + 1$: 61, 121, 181, 241, 301. Yang habis dibagi 7
          adalah 301.}
    \item Ada berapa pasang bilangan $(a, b)$ dengan $a < b$ yang FPB-nya 12 dan KPK-nya 360?
          \jawab{4}{Tulis $a = 12x$ dan $b = 12y$, dengan $x \times y = 30$ dan $x, y$ tidak punya
          faktor sama. Pasangannya (1, 30), (2, 15), (3, 10), (5, 6).}
    \item Tunjukkan bahwa FPB$(n, n+2)$ hanya bisa 1 atau 2. Kapan nilainya 2?
          \jawab{bukti}{FPB membagi selisihnya, yaitu 2. Nilainya 2 jika $n$ genap.}
    \item Tentukan FPB$(2^{20} - 1, 2^{12} - 1)$.
          \textit{(Petunjuk: $2^{20} - 1 = 2^8 \times (2^{12} - 1) + (2^8 - 1)$.)}
          \jawab{15}{Seperti Langkah 4, FPB-nya sama dengan FPB$(2^{12} - 1, 2^8 - 1)$, lalu
          FPB$(2^8 - 1, 2^4 - 1)$. Karena $255 = 15 \times 17$, hasilnya $2^4 - 1 = 15$.}
\end{enumerate}

\level{Level 4: Setara OSP (Provinsi)}
\textit{Tulis jawaban lengkap dengan alasannya.}
\begin{enumerate}
    \setcounter{enumi}{15}
    \item Buktikan bahwa untuk setiap bilangan asli $n$, pecahan $\dfrac{21n + 4}{14n + 3}$
          tidak dapat disederhanakan.
          \jawab{bukti}{Pecahan tidak dapat disederhanakan jika FPB pembilang dan penyebutnya 1.
          Misalkan $d$ membagi $21n + 4$ dan $14n + 3$. Maka $d$ juga membagi
          \[ 3(14n + 3) - 2(21n + 4) = 42n + 9 - 42n - 8 = 1. \]
          Bilangan yang membagi 1 hanya 1. Jadi FPB-nya 1, dan pecahan itu tidak dapat
          disederhanakan. \hfill$\blacksquare$\\
          \emph{Catatan: ini soal IMO 1959 nomor 1.}}
\end{enumerate}

\level{Level 5: Setara OSN (Nasional)}
\textit{Tulis jawaban lengkap dengan alasannya.}
\begin{enumerate}
    \setcounter{enumi}{16}
    \item Tentukan semua pasangan bilangan asli $(a, b)$ dengan $a < b$, $a + b = 84$, dan
          FPB$(a, b) = 12$. Buktikan tidak ada pasangan lain.
          \jawab{(12, 72), (24, 60), (36, 48)}{Karena FPB-nya 12, tulis $a = 12x$ dan $b = 12y$
          dengan FPB$(x, y) = 1$ (kalau $x$ dan $y$ punya faktor sama, FPB$(a,b)$ akan lebih dari 12).
          Dari $a + b = 84$ didapat $x + y = 7$ dengan $x < y$. Kemungkinannya $(1, 6)$, $(2, 5)$,
          $(3, 4)$, dan ketiganya memang tidak punya faktor sama selain 1. Jadi pasangannya
          $(12, 72)$, $(24, 60)$, $(36, 48)$. \hfill$\blacksquare$}
\end{enumerate}

\newpage
% ====================================================================
\section*{Soal Campuran}
% ====================================================================

Soal-soal di bawah ini sengaja tidak diberi keterangan topik, sama seperti soal olimpiade
sungguhan. Coba tebak dulu, soal ini pakai cara yang mana? Waktu mengerjakan 60 menit.

\begin{enumerate}
    \item Tentukan angka satuan dari $2026^{2026} + 2027^{2027}$.
          \jawab{9}{$6 + 3$.}
    \item Tentukan sisa pembagian $2^{2026}$ oleh 5.
          \jawab{4}{Angka satuan $2^{2026}$ adalah 4.}
    \item Ada berapa bilangan $n$ dari 1 sampai 100 sehingga $n^2 + n$ genap?
          \jawab{100}{$n^2 + n = n(n+1)$ selalu genap.}
    \item Tentukan bilangan 4 digit terbesar yang habis dibagi 12 dan keempat digitnya berbeda.
          \jawab{9876}{76 habis dibagi 4 dan jumlah digitnya 30.}
    \item Berapakah banyak faktor dari $10!$?
          \jawab{270}{$10! = 2^8 \times 3^4 \times 5^2 \times 7$, jadi $9 \times 5 \times 3 \times 2$.}
    \item Hari ini hari Rabu. Hari apakah $2^{100}$ hari lagi?
          \jawab{Jumat}{Pola sisa $2^n$ dibagi 7: 2, 4, 1. $100 \div 3$ sisa 1, jadi maju 2 hari.}
    \item Jumlah tiga bilangan prima yang berbeda adalah 40. Tentukan hasil kali ketiganya.
          \jawab{434}{Jumlahnya genap, jadi salah satunya harus 2. Dua lainnya berjumlah 38, yaitu 7 dan 31.}
    \item Berapakah bilangan terkecil $n$ sehingga $2026 \times n$ adalah bilangan kuadrat?
          \jawab{2026}{$2026 = 2 \times 1013$, dan 1013 adalah bilangan prima.}
    \item Tentukan bilangan terkecil yang habis dibagi 1, 2, 3, \dots, 10.
          \jawab{2520}{$2^3 \times 3^2 \times 5 \times 7$.}
    \item Tentukan sisa pembagian $1! + 2! + 3! + \dots + 2026!$ oleh 12.
          \jawab{9}{Mulai $4! = 24$, semuanya habis dibagi 12. Sisanya $1 + 2 + 6$.}
\end{enumerate}

\end{document}
